\section{2019 Analysis \& Differential Equations}

\subsection{Individual}

\begin{exercise}
Let $F:\mathbb{R}\to\mathbb{R}$ be a strictly convex function. Let $u:[0,1]\to\mathbb{R}$ be a continuous function, with
\[
\int_0^1 u(x)\,dx = 0.
\]
Show that
\[
\int_0^1 F(u(x))\,dx \leq \frac{F(\|u\|_{\infty}) + F(-\|u\|_{\infty})}{2}
\]
where $\|u\|_{\infty} := \sup_{x\in[0,1]}|u(x)|$. Also determine when equality occurs.
\end{exercise}

\begin{solution}
Let $M = \|u\|_{\infty}$. Then for all $x \in [0, 1]$, $-M \leq u(x) \leq M$. Since $F$ is a convex function, for any $y \in [-M, M]$, we can express $y$ as a convex combination of $-M$ and $M$:
\[
y = \lambda M + (1-\lambda)(-M), \quad \text{where } \lambda = \frac{y+M}{2M}.
\]
By the definition of convexity:
\[
F(y) \leq \frac{y+M}{2M} F(M) + \frac{M-y}{2M} F(-M).
\]
Substituting $y = u(x)$ and integrating over $[0, 1]$:
\[
\int_0^1 F(u(x))\,dx \leq \int_0^1 \left( \frac{u(x)+M}{2M} F(M) + \frac{M-u(x)}{2M} F(-M) \right) dx.
\]
Using the linearity of the integral and the fact that $\int_0^1 u(x)\,dx = 0$:
\[
\int_0^1 F(u(x))\,dx \leq \frac{F(M)}{2M} \int_0^1 (u(x)+M) dx + \frac{F(-M)}{2M} \int_0^1 (M-u(x)) dx
\]
\[
= \frac{F(M)}{2M} (0 + M) + \frac{F(-M)}{2M} (M - 0) = \frac{F(M) + F(-M)}{2}.
\]
This proves the inequality.

Now, we determine when equality occurs. Since $F$ is strictly convex, the inequality $F(y) \leq \lambda F(M) + (1-\lambda)F(-M)$ is strict for all $y \in (-M, M)$. Equality in the integral occurs if and only if $u(x) \in \{M, -M\}$ for almost all $x \in [0, 1]$. However, $u$ is given as a continuous function. If $M > 0$, a continuous function that takes values in $\{M, -M\}$ must be constant by the Intermediate Value Theorem. Thus, $u \equiv M$ or $u \equiv -M$. But the condition $\int_0^1 u(x)\,dx = 0$ would then imply $M = 0$, which means $u \equiv 0$. If $M=0$, then $u \equiv 0$ and equality clearly holds. Thus, equality occurs if and only if $u \equiv 0$.
\end{solution}

\begin{exercise}
Prove that there exists a universal constant $K$, for all $C^1$ function $f:\mathbb{R}^2\to\mathbb{R}$, if $f\in L^1(\mathbb{R}^2)\cap L^2(\mathbb{R}^2)$ and $|\nabla f|\in L^2(\mathbb{R}^2)$, we have the following inequality:
\[
\|f\|_{L^2(\mathbb{R}^2)}^2 \leq K\|f\|_{L^1(\mathbb{R}^2)}\|\nabla f\|_{L^2(\mathbb{R}^2)}.
\]

Can you provide constant $K$ so that $K<10$? In the problem, all the $L^p$-spaces are defined with respect to the Lebesgue measure.
\end{exercise}

\begin{solution}
By the Plancherel theorem, $\|f\|_{L^2}^2 = \|\hat{f}\|_{L^2}^2$, where $\hat{f}$ is the Fourier transform. We split the integral of $|\hat{f}|^2$ into two regions: $|\xi| \leq R$ and $|\xi| > R$.
\[
\int_{\mathbb{R}^2} |\hat{f}(\xi)|^2 d\xi = \int_{|\xi| \leq R} |\hat{f}(\xi)|^2 d\xi + \int_{|\xi| > R} \frac{1}{|\xi|^2} |\xi|^2 |\hat{f}(\xi)|^2 d\xi.
\]
We know that $\|\hat{f}\|_{L^\infty} \leq \|f\|_{L^1}$ and $\int |\xi|^2 |\hat{f}(\xi)|^2 d\xi = \frac{1}{4\pi^2} \|\nabla f\|_{L^2}^2$ (depending on the normalization of the Fourier transform). Using the standard normalization $\hat{f}(\xi) = \int f(x) e^{-2\pi i x \cdot \xi} dx$:
\[
\int_{|\xi| \leq R} |\hat{f}(\xi)|^2 d\xi \leq \|\hat{f}\|_{L^\infty}^2 \cdot \text{Area}(\{|\xi| \leq R\}) = \|f\|_{L^1}^2 \pi R^2.
\]
For the second term:
\[
\int_{|\xi| > R} \frac{1}{|\xi|^2} |\xi|^2 |\hat{f}(\xi)|^2 d\xi \leq \frac{1}{R^2} \int_{\mathbb{R}^2} |\xi|^2 |\hat{f}(\xi)|^2 d\xi = \frac{1}{R^2} \frac{1}{4\pi^2} \|\nabla f\|_{L^2}^2.
\]
Thus, $\|f\|_{L^2}^2 \leq \pi R^2 \|f\|_{L^1}^2 + \frac{1}{4\pi^2 R^2} \|\nabla f\|_{L^2}^2$. Minimizing with respect to $R^2$:
Setting $A = \pi \|f\|_{L^1}^2$ and $B = \frac{1}{4\pi^2} \|\nabla f\|_{L^2}^2$, the minimum occurs at $R^2 = \sqrt{B/A}$.
The minimum value is $2\sqrt{AB} = 2 \sqrt{\pi \|f\|_{L^1}^2 \frac{1}{4\pi^2} \|\nabla f\|_{L^2}^2} = \frac{1}{\sqrt{\pi}} \|f\|_{L^1} \|\nabla f\|_{L^2}$.
Thus, $K = \frac{1}{\sqrt{\pi}} \approx 0.564$, which is well below 10.
\end{solution}

\begin{exercise}
Let $f:\mathbb{R}^2\to\mathbb{R}$ be a harmonic function. Suppose
\[
\lim_{|x|\to\infty}\frac{|f(x)|}{\ln|x|} = 0.
\]
Prove or disprove that $f$ is a constant.
\end{exercise}

\begin{solution}
This is a variation of Liouville's Theorem for harmonic functions. A harmonic function $f$ on $\mathbb{R}^2$ can be viewed as the real part of an entire function $g(z)$. If $f$ is harmonic and satisfies $|f(z)| \leq C(1+|z|^k)$, then $f$ is a polynomial of degree at most $k$.
In our case, the condition $\lim_{|x|\to\infty} \frac{|f(x)|}{\ln |x|} = 0$ implies that for any $\epsilon > 0$, there exists $R$ such that for $|x| > R$, $|f(x)| \leq \epsilon \ln |x|$. Since $\ln |x| < |x|^\alpha$ for any $\alpha > 0$ and sufficiently large $|x|$, we have $|f(x)| \leq C_\alpha (1+|x|^\alpha)$ for all $\alpha > 0$.
Taking $\alpha < 1$, the only polynomials of degree less than 1 are constants. Thus $f$ must be a constant.
Alternatively, one can use the mean value property and gradient estimates for harmonic functions. For any $x_0 \in \mathbb{R}^2$ and $R > 0$:
\[
|\nabla f(x_0)| \leq \frac{2}{R} \sup_{B_R(x_0)} |f|.
\]
With the given growth condition, as $R \to \infty$, the right hand side goes to zero because $\frac{\ln R}{R} \to 0$. Thus $\nabla f \equiv 0$, and $f$ is constant.
\end{solution}

\begin{exercise}
\begin{itemize}
\item[(a)] Show that there does not exist a holomorphic function $f$ on $\mathbb{C}\setminus\{1,-1\}$ so that
\[
f'(z) = \frac{1}{(z^2-1)^{2019}} \quad \text{for all } z\in\mathbb{C}\setminus\{1,-1\}.
\]

\item[(b)] Show that there exist a set $L\subset\mathbb{C}$ and a holomorphic function $F^{\dagger}$ on $\mathbb{C}\setminus L$ so that $L$ has Hausdorff dimension $1$, and
\[
F^{\dagger}(z) = \int_{z_0}^z \frac{1}{(w^2-1)^{2019}} dw \quad \text{is well-defined}.
\]
\end{itemize}
\end{exercise}

\begin{solution}
(a) A holomorphic function $f$ with a given derivative $g(z) = (z^2-1)^{-2019}$ exists on a domain $D$ if and only if the residue of $g$ at every singularity inside any loop in $D$ is zero (or more generally, the integral of $g$ along any closed path is zero).
Let $g(z) = \frac{1}{(z-1)^{2019}(z+1)^{2019}}$. The residues at $z=1$ and $z=-1$ can be calculated using the formula for poles of order $n$:
\[
\text{Res}(g, 1) = \frac{1}{2018!} \frac{d^{2018}}{dz^{2018}} \left( (z+1)^{-2019} \right) \Big|_{z=1}.
\]
The $k$-th derivative of $(z+1)^{-2019}$ is $(-2019)(-2020)\dots(-2019-k+1)(z+1)^{-2019-k}$.
For $k=2018$, this is non-zero. Specifically, $\text{Res}(g, 1) = \frac{(-1)^{2018} (2019+2018-1)!}{(2019-1)! 2018! (1+1)^{2019+2018}} = \frac{4036!}{2018! 2018! 2^{4037}} \neq 0$.
Since the residue is non-zero, the integral along a small circle around $z=1$ is non-zero, so no global primitive $f$ exists on $\mathbb{C} \setminus \{1, -1\}$.

(b) If we choose $L$ to be a simple path (e.g., a line segment) connecting $-1$ and $1$, then $D = \mathbb{C} \setminus L$ is simply connected. In any simply connected domain, every holomorphic function has a primitive. Thus, there exists a holomorphic function $F^\dagger$ on $\mathbb{C} \setminus L$ such that $(F^\dagger)'(z) = (z^2-1)^{-2019}$. The set $L$ (a line segment) has Hausdorff dimension 1.
\end{solution}

\begin{exercise}
Let $\Omega\subset\mathbb{R}^2$ be a bounded domain with smooth boundary. Prove that, for all $p>1$ and $1\leq q < \infty$, for all $f\in L^p(\Omega)$, there exists a unique $u\in H_0^1(\Omega)$, such that
\[
\triangle u = |u|^{q-1}u + f \text{ in } \Omega.
\]
\end{exercise}

\begin{solution}
Consider the energy functional $J: H_0^1(\Omega) \to \mathbb{R}$ defined by:
\[
J(u) = \int_\Omega \left( \frac{1}{2}|\nabla u|^2 + \frac{1}{q+1}|u|^{q+1} + fu \right) dx.
\]
Since $\Omega$ is bounded in $\mathbb{R}^2$, by the Sobolev embedding theorem, $H_0^1(\Omega) \subset L^s(\Omega)$ for all $1 \leq s < \infty$. Thus the term $\int |u|^{q+1}$ is well-defined for all $q < \infty$.
The functional $J$ is strictly convex because $u \mapsto \frac{1}{2}|\nabla u|^2$ is strictly convex on $H_0^1(\Omega)$ and $u \mapsto \frac{1}{q+1}|u|^{q+1}$ is convex for $q \geq 1$.
Coercivity: Using Poincaré's inequality and Hölder's inequality:
\[
J(u) \geq \frac{1}{2}\|\nabla u\|_2^2 + \frac{1}{q+1}\|u\|_{q+1}^{q+1} - \|f\|_{p} \|u\|_{p'}.
\]
Since $p > 1$, $p' < \infty$, and $\|u\|_{p'} \leq C \|\nabla u\|_2$. For large $\|\nabla u\|_2$, the quadratic term dominates the linear term, so $J(u) \to \infty$ as $\|u\|_{H_0^1} \to \infty$.
By the direct method in the calculus of variations, $J$ attains a global minimum $u \in H_0^1(\Omega)$. The Euler-Lagrange equation for $J$ is exactly $-\triangle u + |u|^{q-1}u + f = 0$, which is $\triangle u = |u|^{q-1}u + f$. Uniqueness follows from the strict convexity of $J$.
\end{solution}

\subsection{Team}

\begin{exercise}
Show that there is no non-zero $f\in C_0^{\infty}(\mathbb{R}^2)$ (compactly supported smooth function) so that its Fourier transform $\widehat{f}(\xi)$ is also compactly supported.
\end{exercise}

\begin{solution}
Suppose $f \in C_0^\infty(\mathbb{R}^2)$ has compact support. By the Paley-Wiener theorem, its Fourier transform $\hat{f}(\xi)$ can be extended to an entire function on $\mathbb{C}^2$. If $\hat{f}$ is also compactly supported on $\mathbb{R}^2$, then $\hat{f}$ must vanish on an open set in $\mathbb{R}^2$. By the identity theorem for holomorphic functions, an entire function that vanishes on an open subset of $\mathbb{R}^2$ must be identically zero. Therefore, $\hat{f} \equiv 0$, which implies $f \equiv 0$ by the Fourier inversion formula. Thus, no such non-zero function exists.
\end{solution}

\begin{exercise}
Prove the following classical interior Schauder estimates:

There exists a universal constant $C$, for all smooth compactly supported functions $u,f\in C_0^{\infty}(\mathbb{R}^3)$ with $\Delta u = f$, we have
\[
\|u\|_{C^{2,\alpha}} \leq C\|f\|_{C^{0,\alpha}},
\]
where $0<\alpha<1$ and $\|\cdot\|_{C^{k,\alpha}}$ are H\"older norms.
\end{exercise}

\begin{solution}
The solution to $\Delta u = f$ in $\mathbb{R}^3$ is given by the Newtonian potential:
\[
u(x) = \int_{\mathbb{R}^3} \Phi(x-y) f(y) dy, \quad \text{where } \Phi(x) = -\frac{1}{4\pi|x|}.
\]
The second derivatives are given by $D^2 u(x) = PV \int_{\mathbb{R}^3} D^2 \Phi(x-y) f(y) dy + \frac{1}{3} f(x) I$.
This is a singular integral of the form $Tf(x) = \lim_{\epsilon \to 0} \int_{|x-y|>\epsilon} K(x-y) f(y) dy$. The kernel $K(z) = D^2 \Phi(z)$ is homogeneous of degree $-3$ and has zero mean on the unit sphere.
The Schauder estimates state that such singular integral operators are bounded from $C^{0,\alpha}$ to $C^{0,\alpha}$ for $0 < \alpha < 1$.
Specifically, one can show that $|D^2 u(x) - D^2 u(y)| \leq C \|f\|_{C^{0,\alpha}} |x-y|^\alpha$ by splitting the integral into regions $|x-z| < 2|x-y|$ and $|x-z| \geq 2|x-y|$, utilizing the cancellation property of the kernel $K$. Since $u$ is also bounded in $L^\infty$ and its first derivatives are bounded (by potential estimates), we obtain $\|u\|_{C^{2,\alpha}} \leq C \|f\|_{C^{0,\alpha}}$.
\end{solution}

\begin{exercise}
Let $(X,\mathcal{A},\mu)$ be a probability space and let $T:X\to X$ be a measurable and measure preserving map, i.e., for all $A\in\mathcal{A}$, we have $\mu(T^{-1}(A))=\mu(A)$. For $A,B\in\mathcal{A}$, if $\mu(A-B)=\mu(B-A)=0$, we say that $A=B$ a.e.

Assume $A\in\mathcal{A}$ such that $T^{-1}(A)=A$ a.e. Prove that there exists a set $B\in\mathcal{A}$ so that $T^{-1}B=B$ and $A=B$ a.e.
\end{exercise}

\begin{solution}
Let $A \in \mathcal{A}$ such that $\mu(T^{-1}(A) \triangle A) = 0$.
Define $B = \limsup_{n \to \infty} T^{-n}(A) = \bigcap_{k=1}^\infty \bigcup_{n=k}^\infty T^{-n}(A)$.
First, we check if $T^{-1}B = B$.
$T^{-1}B = T^{-1} \left( \bigcap_{k=1}^\infty \bigcup_{n=k}^\infty T^{-n}(A) \right) = \bigcap_{k=1}^\infty \bigcup_{n=k}^\infty T^{-(n+1)}(A) = \bigcap_{k=1}^\infty \bigcup_{j=k+1}^\infty T^{-j}(A) = B$.
Now we show $A = B$ a.e.
Since $\mu(T^{-1}(A) \triangle A) = 0$ and $T$ is measure preserving, we have $\mu(T^{-(n+1)}(A) \triangle T^{-n}(A)) = 0$ for all $n \geq 0$.
By induction, $\mu(T^{-n}(A) \triangle A) = 0$ for all $n$.
Then $\mu( \bigcup_{n=k}^\infty T^{-n}(A) \triangle A ) \leq \sum_{n=k}^\infty \mu(T^{-n}(A) \triangle A) = 0$ is not directly applicable since the sum might diverge.
However, $A = T^{-n}(A)$ a.e. implies $\chi_A = \chi_{T^{-n}(A)}$ a.e.
Then $\sup_{n \geq k} \chi_{T^{-n}(A)} = \chi_A$ a.e. for any $k$.
Thus $\chi_B = \lim_{k \to \infty} \sup_{n \geq k} \chi_{T^{-n}(A)} = \chi_A$ a.e.
So $B = A$ a.e.
\end{solution}

\begin{exercise}
Is there an entire function $f$ with infinitely many zeroes, so that for every $r\in(0,1)$, there exist constants $A_r,B_r<\infty$ such that
\[
|f(z)| \leq A_r e^{B_r|z|^r}
\]
for every $z\in\mathbb{C}$?
\end{exercise}

\begin{solution}
Yes, such functions exist. These are entire functions of order 0. The order $\rho$ of an entire function $f$ is defined as $\rho = \limsup_{R \to \infty} \frac{\ln \ln M(R)}{\ln R}$, where $M(R) = \max_{|z|=R} |f(z)|$. The given condition implies $\rho \leq r$ for all $r \in (0, 1)$, hence $\rho = 0$.
By the Weierstrass Factorization Theorem, we can construct such a function using its zeros $\{a_n\}$. For a function of order 0, it is sufficient that the zeros grow fast enough. Let $a_n = e^n$. Then the number of zeros in a disk of radius $R$ is $n(R) \approx \ln R$.
The canonical product $f(z) = \prod_{n=1}^\infty (1 - z/a_n)$ defines an entire function. The order of this product is equal to the exponent of convergence of its zeros, $\lambda = \inf \{ \alpha > 0 : \sum |a_n|^{-\alpha} < \infty \}$.
For $a_n = e^n$, $\sum (e^n)^{-\alpha}$ converges for all $\alpha > 0$. Thus $\lambda = 0$, and the function $f(z)$ has order 0, satisfying the required growth condition.
\end{solution}

\begin{exercise}
Let $u(t,x,y)$ be a smooth real function defined on $\mathbb{R}\times\mathbb{R}^2$ where $t\in\mathbb{R}$ and $(x,y)\in\mathbb{R}^2$. We assume that it solves the following semilinear wave equation:
\[
-\frac{\partial^2}{\partial t^2}u + \triangle u = u^3.
\]
If the supports of the initial data $u(0,x)$ and $\frac{\partial u}{\partial t}(0,x)$ are compact, prove that, for all $t_0\in\mathbb{R}$, the supports of $u(t_0,x)$ and $\frac{\partial u}{\partial t}(t_0,x)$ are compact.
\end{exercise}

\begin{solution}
This follows from the finite speed of propagation property of the wave equation. Let $K \subset \mathbb{R}^2$ be a compact set containing the supports of $u(0,x)$ and $u_t(0,x)$.
Consider a point $x_0$ such that $\text{dist}(x_0, K) > |t_0|$. The backward light cone $C = \{ (t, x) : |x - x_0| \leq |t_0| - t, 0 \leq t \leq |t_0| \}$ (assuming $t_0 > 0$) has its base in $\mathbb{R}^2 \setminus K$.
We define the local energy $E(t) = \int_{|x-x_0| \leq |t_0|-t} \left( \frac{1}{2} u_t^2 + \frac{1}{2} |\nabla u|^2 + \frac{1}{4} u^4 \right) dx$.
Differentiating $E(t)$ and using the equation $u_{tt} - \triangle u = -u^3$:
\[
\frac{dE}{dt} = \int_{|x-x_0| \leq |t_0|-t} (u_t u_{tt} + \nabla u \cdot \nabla u_t + u^3 u_t) dx - \int_{|x-x_0| = |t_0|-t} \left( \frac{1}{2} u_t^2 + \frac{1}{2} |\nabla u|^2 + \frac{1}{4} u^4 \right) d\sigma
\]
\[
= \int_{|x-x_0| \leq |t_0|-t} \text{div}(u_t \nabla u) dx - \int_{|x-x_0| = |t_0|-t} (\dots) d\sigma = \int_{|x-x_0| = |t_0|-t} (u_t \frac{\partial u}{\partial n} - \frac{1}{2} u_t^2 - \frac{1}{2} |\nabla u|^2 - \frac{1}{4} u^4) d\sigma.
\]
Since $|u_t \frac{\partial u}{\partial n}| \leq |u_t| |\nabla u| \leq \frac{1}{2} u_t^2 + \frac{1}{2} |\nabla u|^2$, we have $\frac{dE}{dt} \leq 0$.
Since $E(0) = 0$ (because the support of initial data is outside the base of the cone), it follows that $E(t) = 0$ for all $0 \leq t \leq |t_0|$.
In particular, $E(|t_0|) = 0$, which implies $u(t_0, x_0) = 0$ and $u_t(t_0, x_0) = 0$.
Thus the support of $u$ at time $t_0$ is contained in the set $\{ x : \text{dist}(x, K) \leq |t_0| \}$, which is compact.
\end{solution}


